\chapter{Le modèle binomial}\label{ch:dv:the-binomial-model}

Le tableau blanc montre une action à 100 et trois pas de temps. À chaque pas, le
cours de l'action monte de 10~\% ou baisse de 10~\%, et un dépôt bancaire croît de 2~\%. La question est le
prix d'un call de strike 100 qui expire après le troisième pas. Il n'existe
aucun nombre de volatilité au tableau, aucune probabilité de hausse et aucune formule,
et la réponse, 10.3603, n'en a besoin d'aucun. Chaque prix de ce chapitre est le
coût d'un portefeuille d'actions et de dépôt qui verse exactement ce que l'option
rapporte, nœud par nœud~; la probabilité qui apparaît en chemin est celle
que le chapitre~1 appelle risque-neutre, et c'est un résultat des prix, non une donnée.
L'arbre est le plus petit \omterm{def:dv:no-arbitrage-and-the-fundamental-theorems:complete}{marché complet} dans lequel les options ne sont pas triviales,
et c'est aussi un outil de pricing opérationnel~: avec assez de pas, il converge vers la
valeur de Black--Scholes du chapitre~3, et il price l'exercice américain, ce
qu'aucune formule fermée ne permet.

\section{Une période~: réplication et probabilité risque-neutre}

\begin{definition}[Modèle binomial]\label{def:dv:the-binomial-model:binomial}
Dans le \emph{modèle binomial}\index{modèle binomial} le temps se découpe en $n$ pas de
longueur $\Delta t=T/n$. À chaque pas, le cours de l'action est multiplié par $u$ ou par
$d<u$, et une unité de cash déposée croît jusqu'à $R=e^{r\Delta t}$ (avec un
rendement de dividende continu $q$, la croissance propre de l'action se compare à
$e^{(r-q)\Delta t}$). Le modèle est sans arbitrage si et seulement si
$d<e^{(r-q)\Delta t}<u$.
\end{definition}

Sur un pas, un produit versant $V_u$ après un mouvement haussier et $V_d$ après un
mouvement baissier est répliqué par $\Delta$ actions et un dépôt $b$ si $\Delta S u+bR=V_u$ et
$\Delta Sd+bR=V_d$. Deux équations, deux inconnues~:
\begin{equation}\label{eq:dv:the-binomial-model:onestep}
\Delta=\frac{V_u-V_d}{S(u-d)},\qquad
V=\Delta S+b=\frac{1}{R}\bigl(pV_u+(1-p)V_d\bigr),\qquad p=\frac{R-d}{u-d}.
\end{equation}
Le poids $p$ appartient à $(0,1)$ exactement sous la condition d'absence d'arbitrage~; c'est
la probabilité risque-neutre du mouvement haussier, le \omterm{def:dv:no-arbitrage-and-the-fundamental-theorems:ad}{prix d'état} normalisé du chapitre~1,
et la probabilité réelle d'un mouvement haussier n'intervient nulle part. Deux
traders en désaccord sur cette probabilité s'accordent sur le prix, car chacun peut
construire l'option à partir d'actions et de cash.

\begin{example}[Un pas]\label{ex:dv:the-binomial-model:onestep}
Une action à 50 vaudra 60 ou 40~; les taux sont nuls. Un call de strike 50 verse 10 ou
0~: $\Delta=10/20=0.5$ action, et le dépôt vaut $b=(0-0.5\times40)=-20$, un
emprunt de 20. Le call coûte $0.5\times50-20=5$, et $p=(1-0.8)/0.4=0.5$.
\end{example}

\section{Plusieurs périodes~: l'arbre recombinant}

\begin{definition}[Arbre recombinant]\label{def:dv:the-binomial-model:recombining}
Un \emph{arbre recombinant}\index{arbre recombinant} est un arbre dans lequel un mouvement haussier
suivi d'un mouvement baissier atteint le même nœud que l'inverse~: après $k$ pas
le cours de l'action se trouve à l'une des $k+1$ valeurs $Su^jd^{k-j}$, $j=0,\dots,k$, au lieu de
$2^k$. L'évaluation sur cet arbre coûte $O(n^2)$ opérations et $O(n)$ en mémoire.
\end{definition}

Appliquer \eqref{eq:dv:the-binomial-model:onestep} nœud par nœud, depuis le dernier
pas jusqu'au premier, revient à l'induction rétrograde de la programmation dynamique
(livre~4, chapitre~9). Le portefeuille de chaque nœud est rééquilibré au nœud suivant
en n'utilisant que sa propre valeur~: c'est une \omterm{def:dv:no-arbitrage-and-the-fundamental-theorems:sf}{stratégie autofinancée} au sens du
chapitre~1, et il réplique le produit sur toute trajectoire.

\begin{proposition}[Le prix binomial]\label{prop:dv:the-binomial-model:formula}
Un produit européen versant $g(S_T)$ après $n$ pas vaut
\[
V_0=R^{-n}\sum_{j=0}^{n}\binom{n}{j}p^j(1-p)^{n-j}\,g\bigl(S_0u^jd^{n-j}\bigr),
\]
l'espérance actualisée du gain lorsque le nombre de mouvements haussiers suit une
loi binomiale de paramètres $n$ et $p$.
\end{proposition}
\begin{proof}
Par récurrence sur le nombre de pas restants~: chaque pas rétrograde remplace deux
valeurs par leur moyenne pondérée par $p$, divisée par $R$, et le nombre de chemins
atteignant le nœud $j$ est $\binom{n}{j}$.
\end{proof}

\begin{omfigure}
\begin{tikzpicture}[font=\footnotesize,x=2.9cm,y=0.95cm,
  nd/.style={draw=omDef,thick,rounded corners=2pt,fill=omDef!6,align=center,inner sep=2pt,minimum width=1.55cm},
  ex/.style={nd,draw=omThm,fill=omThm!8}]
\node[nd] (a) at (0,0) {100\\4.908};
\node[ex] (b0) at (1,-1.5) {90\\\textbf{10.000}};
\node[nd] (b1) at (1,1.5) {110\\1.676};
\node[ex] (c0) at (2,-3) {81\\\textbf{19.000}};
\node[nd] (c1) at (2,0) {99\\4.275};
\node[nd] (c2) at (2,3) {121\\0.000};
\node[nd] (d0) at (3,-4.5) {72.9\\27.100};
\node[nd] (d1) at (3,-1.5) {89.1\\10.900};
\node[nd] (d2) at (3,1.5) {108.9\\0.000};
\node[nd] (d3) at (3,4.5) {133.1\\0.000};
\foreach \s/\t in {a/b0,a/b1,b0/c0,b0/c1,b1/c1,b1/c2,c0/d0,c0/d1,c1/d1,c1/d2,c2/d2,c2/d3}
  \draw[gray,-{Stealth}] (\s) -- (\t);
\node[anchor=west,align=left,text=omThm] at (0.1,-3.9) {rouge~: exercice immédiat\\supérieur à la détention};
\end{tikzpicture}

\omcaption{Arbre du tableau blanc~: cours de l'action (ligne du haut) et valeur du put américain
(ligne du bas) à chaque nœud, $u=1.1$, $d=0.9$, 2~\% par pas, strike 100. À 90
et à 81, le put vaut plus exercé (10 et 19) que conservé (9.196 et
17.039), et l'arbre retient la plus grande valeur. Données~: le tutoriel.}\label{fig:dv:the-binomial-model:tree}
\end{omfigure}

\begin{example}[Le call du tableau blanc]\label{ex:dv:the-binomial-model:whiteboard}
Avec $u=1.1$, $d=0.9$, $R=1.02$~: $p=(1.02-0.9)/0.2=0.6$. Après trois pas, le
cours de l'action vaut 133.1, 108.9, 89.1 ou 72.9, et le call verse 33.1, 8.9, 0 ou 0. Sa
valeur est $(0.6^3\times33.1+3\times0.6^2\times0.4\times8.9)/1.02^3=10.3603$. Au
premier nœud, le \omterm{def:dv:no-arbitrage-and-the-fundamental-theorems:claim}{portefeuille de réplication} détient $\Delta=0.6240$ action et un
dépôt de $-52.04$~; la parité call-put tient sur l'arbre, $C-P=100-100/1.02^3=5.7678$.
\end{example}

\section{Calibrer l'arbre à une volatilité}

Pour modéliser une action de volatilité $\sigma$, les facteurs multiplicatifs du pas doivent reproduire, au premier
ordre en $\Delta t$, la moyenne et la variance du log-rendement sur un pas. Deux
choix sont classiques~; un troisième est conçu pour gagner en précision.

\begin{definition}[Arbres de Cox--Ross--Rubinstein et de Jarrow--Rudd]\label{def:dv:the-binomial-model:crr}
L'\emph{arbre de Cox--Ross--Rubinstein}\index{arbre de Cox--Ross--Rubinstein} (Cox--Ross--Rubinstein tree) prend
$u=e^{\sigma\sqrt{\Delta t}}$ et $d=1/u$, de sorte que l'arbre soit symétrique en
log-prix et la dérive n'intervient que par $p$. L'\emph{arbre de
Jarrow--Rudd}\index{arbre de Jarrow--Rudd} (Jarrow--Rudd tree) prend
$u,d=\exp\bigl((r-q-\tfrac12\sigma^2)\Delta t\pm\sigma\sqrt{\Delta t}\bigr)$, ce qui
fait passer la dérive dans les facteurs et rend $p$ proche d'un demi.
\end{definition}

\begin{proposition}[Appariement des moments]\label{prop:dv:the-binomial-model:moments}
Dans les deux arbres, le log-rendement sur un pas a une moyenne neutre au risque
$(r-q-\tfrac12\sigma^2)\Delta t+O(\Delta t^{3/2})$ et une variance
$\sigma^2\Delta t+O(\Delta t^2)$, et la croissance de l'action coïncide avec
$e^{(r-q)\Delta t}$ exactement, grâce au choix de $p$.
\end{proposition}
\begin{proof}
Pour l'\omterm{def:dv:the-binomial-model:crr}{arbre de Cox--Ross--Rubinstein}, on développe la probabilité
$p=(e^{(r-q)\Delta t}-e^{-\sigma\sqrt{\Delta t}})/(2\sinh\sigma\sqrt{\Delta t})=
\tfrac12+\tfrac{(r-q-\sigma^2/2)\sqrt{\Delta t}}{2\sigma}+O(\Delta t)$~; la moyenne de
$\pm\sigma\sqrt{\Delta t}$ vaut $(2p-1)\sigma\sqrt{\Delta t}$, et son moment d'ordre 2
vaut $\sigma^2\Delta t$. Le calcul pour l'\omterm{def:dv:the-binomial-model:crr}{arbre de Jarrow--Rudd} est le même, à ceci près que la dérive
est reportée dans les facteurs.
\end{proof}

Par le théorème central limite (livre~4, chapitre~1), le nombre de mouvements à la hausse,
centré et réduit, devient gaussien~: le log-prix terminal de l'arbre converge
en loi vers la loi normale du mouvement brownien géométrique
(voir la \cref{fig:dv:the-binomial-model:terminal}), et le prix de l'arbre d'une option européenne
converge vers le prix de Black--Scholes du chapitre~3.

\begin{omfigure}
\begin{tikzpicture}
\begin{axis}[width=10.5cm,height=4.8cm,ybar,bar width=4pt,
  xlabel={log-rendement à l'échéance},ylabel={densité},
  tick label style={font=\small},label style={font=\small},
  xmin=-0.85,xmax=0.85,ymin=0,ymax=2.4,ymajorgrids,grid style={gray!25},
  legend columns=2,legend style={font=\footnotesize,at={(0.5,-0.3)},anchor=north,draw=none,fill=none,/tikz/every even column/.append style={column sep=8pt}},
  table/col sep=comma]
\addplot[fill=omDef!55,draw=omDef] table[x=x,y=tree]{figdata/derivatives/02-the-binomial-model/terminal.csv};
\addplot[omThm,very thick,sharp plot,mark=none] table[x=x,y=normal]{figdata/derivatives/02-the-binomial-model/terminal.csv};
\legend{arbre à 50 pas (probabilité / écart),limite normale}
\end{axis}
\end{tikzpicture}

\omcaption{Loi neutre au risque du log-rendement terminal dans un arbre à 50 pas de
Cox--Ross--Rubinstein (un an, $r=5\%$, $\sigma=20\%$), chaque probabilité
divisée par l'écart $2\sigma\sqrt{\Delta t}$ entre nœuds, comparée à la densité
normale de moyenne $(r-\sigma^2/2)T$ et d'écart type $\sigma\sqrt T$. Données~:
le tutoriel.}\label{fig:dv:the-binomial-model:terminal}
\end{omfigure}

\section{Convergence et ses oscillations}

Pour une option de vente à la monnaie d'un an avec $r=5\%$ et $\sigma=20\%$, la valeur de Black--Scholes
vaut 5.5735. Le prix de Cox--Ross--Rubinstein à 100 pas est inférieur de 0.0200
à celui-ci, à 101 pas il est supérieur de 0.0174~: l'erreur décroît en $1/n$ et change de signe
selon la parité de $n$ (voir la \cref{fig:dv:the-binomial-model:conv}). La cause est le
coude du payoff. Les nœuds proches du prix d'exercice se déplacent par rapport à lui quand $n$ varie,
et le payoff linéaire par morceaux est évalué en un point différent de son coude
à chaque fois~; avec le prix d'exercice exactement égal au spot, un $n$ pair place un nœud sur le coude
et un $n$ impair l'encadre.

\begin{definition}[Arbre de Leisen--Reimer]\label{def:dv:the-binomial-model:lr}
L'\emph{arbre de Leisen--Reimer}\index{arbre de Leisen--Reimer} (Leisen--Reimer tree) utilise un nombre impair de
pas et choisit $p$ et $u$ de sorte que les probabilités de terminer au-dessus
du prix d'exercice, sous la mesure neutre au risque et sous celle de l'action, soient égales à $\Phi(d_2)$ et
$\Phi(d_1)$ (par une inversion normale-binomiale due à Peizer et Pratt).
L'arbre est centré sur le prix d'exercice, et l'erreur pour les options européennes décroît
en $1/n^2$ sans oscillation.
\end{definition}

\begin{omfigure}
\begin{tikzpicture}
\begin{axis}[width=11cm,height=5cm,
  xlabel={nombre de pas $n$},ylabel={prix de l'arbre $-$ formule},
  tick label style={font=\small},label style={font=\small},
  xmin=10,xmax=200,ymin=-0.1,ymax=0.1,grid=major,grid style={gray!25},
  yticklabel style={/pgf/number format/.cd,fixed,precision=2},
  legend columns=2,legend style={font=\footnotesize,at={(0.5,-0.3)},anchor=north,draw=none,fill=none,/tikz/every even column/.append style={column sep=8pt}},
  table/col sep=comma]
\addplot[omDef,thick] table[x=n,y=crr]{figdata/derivatives/02-the-binomial-model/conv_even.csv};
\addplot[omDef,thick,dashed] table[x=n,y=crr]{figdata/derivatives/02-the-binomial-model/conv_odd.csv};
\addplot[omProp,thick] table[x=n,y=jr]{figdata/derivatives/02-the-binomial-model/conv_even.csv};
\addplot[omThm,very thick] table[x=n,y=lr]{figdata/derivatives/02-the-binomial-model/conv_lr.csv};
\legend{Cox--Ross--Rubinstein $n$ pair,Cox--Ross--Rubinstein $n$ impair,Jarrow--Rudd $n$ pair,Leisen--Reimer $n$ impair}
\end{axis}
\end{tikzpicture}

\omcaption{Erreur de trois arbres pour l'option de vente européenne à la monnaie d'un an
($r=5\%$, $\sigma=20\%$, formule 5.5735), nombres de pas pairs et impairs tracés
séparément. L'erreur de Cox--Ross--Rubinstein est négative pour $n$ pair, positive
pour $n$ impair, et décroît en $1/n$~; quant à
l'erreur de Leisen--Reimer est invisible à cette échelle dès 11 pas. Données~: le
tutoriel.}\label{fig:dv:the-binomial-model:conv}
\end{omfigure}

\begin{remark}[Coût de l'oscillation]\label{rem:dv:the-binomial-model:cost}
L'option de vente de Cox--Ross--Rubinstein ne reste à moins d'un cent de la formule qu'à partir de 199
pas~; l'\omterm{def:dv:the-binomial-model:lr}{arbre de Leisen--Reimer} reste à moins de 0.3 cent à 11 pas, et à moins de
$10^{-4}$ à 101 pas. La moyenne des prix en $n$ et $n+1$ supprime l'essentiel de
l'oscillation, au prix d'un second arbre~; l'extrapolation de Richardson (livre~4,
chapitre~27) supprime le terme en $1/n$ dès que l'oscillation a disparu. Le chapitre
22 compare les arbres aux grilles de différences finies, en précision et en temps de calcul.
\end{remark}

\section{Exercice américain sur un arbre}

Une option américaine peut être exercée à tout nœud. À chaque nœud, le détenteur compare
la valeur de la poursuite, donnée par \eqref{eq:dv:the-binomial-model:onestep}, à
la valeur de l'exercice immédiat, et l'option vaut la plus grande des deux~:
\[
V_k^j=\max\Bigl(g(S_k^j),\ R^{-1}\bigl(pV_{k+1}^{j+1}+(1-p)V_{k+1}^{j}\bigr)\Bigr).
\]
Le \omterm{def:dv:no-arbitrage-and-the-fundamental-theorems:claim}{portefeuille de réplication} du vendeur reste autofinancé là où le détenteur
poursuit~; là où le détenteur exerce, il est liquidé pour verser la
valeur d'exercice. Si un détenteur omet d'exercer à un nœud où l'exercice est optimal, le
portefeuille du vendeur vaut plus que l'option à partir de là~: le vendeur
tire profit de l'erreur du détenteur, un thème du chapitre~6.

\begin{proposition}[Pas d'exercice anticipé d'un call sans dividendes]\label{prop:dv:the-binomial-model:call}
Sur un arbre avec $r\ge0$ et $q=0$, un call américain vaut le call européen~:
l'exercice n'est jamais strictement optimal avant l'échéance.
\end{proposition}
\begin{proof}
À tout nœud, la valeur de poursuite vaut au moins
$R^{-1}\bigl(pS_{k+1}^{j+1}+(1-p)S_{k+1}^{j}\bigr)-KR^{-(n-k)}$, par convexité de la
fonction de paiement et par récurrence, ce qui vaut $S_k^j-KR^{-(n-k)}\ge S_k^j-K$~: le call
vaut plus vivant qu'exercé.
\end{proof}

Un put n'a pas une telle borne~: profondément dans la monnaie, son détenteur préfère
recevoir le strike tout de suite et toucher l'intérêt correspondant (\cref{fig:dv:the-binomial-model:american}).
Sur l'arbre du tableau blanc, le put américain vaut 4.9076, contre 4.5925 pour le
put européen~: une prime d'exercice anticipé de 0.3151, gagnée aux deux nœuds de la
\cref{fig:dv:the-binomial-model:tree} où l'exercice l'emporte.

\begin{omfigure}
\begin{tikzpicture}
\begin{axis}[width=10.5cm,height=5cm,
  xlabel={spot},ylabel={valeur du put},
  tick label style={font=\small},label style={font=\small},
  xmin=60,xmax=140,ymin=0,ymax=42,grid=major,grid style={gray!25},
  legend columns=3,legend style={font=\footnotesize,at={(0.5,-0.3)},anchor=north,draw=none,fill=none,/tikz/every even column/.append style={column sep=8pt}},
  table/col sep=comma]
\addplot[gray,thick,dashed] table[x=spot,y=intrinsic]{figdata/derivatives/02-the-binomial-model/american.csv};
\addplot[omDef,very thick] table[x=spot,y=european]{figdata/derivatives/02-the-binomial-model/american.csv};
\addplot[omThm,very thick] table[x=spot,y=american]{figdata/derivatives/02-the-binomial-model/american.csv};
\legend{valeur d'exercice $(K-S)^+$,européen,américain}
\end{axis}
\end{tikzpicture}

\omcaption{Puts d'un an de strike 100 ($r=5\%$, $\sigma=20\%$, \omterm{def:dv:the-binomial-model:lr}{arbre de
Leisen--Reimer} à 501 pas). Le put européen passe sous sa valeur d'exercice profondément dans la
monnaie, là où attendre coûte l'intérêt sur le strike~; le put américain ne le fait jamais, et
rejoint la valeur d'exercice tangentiellement à la frontière d'exercice, vers
81.1 ici. Données~: le tutoriel.}\label{fig:dv:the-binomial-model:american}
\end{omfigure}

\section{Tutoriel~: un pricer sur arbre et sa convergence}\label{tut:dv:the-binomial-model:tut}

\begin{tutorial}
\textbf{Objectif.} Construire les trois arbres, pricer des européennes et des américaines sur ceux-ci, et
mesurer leur convergence vers la formule. \textbf{État final~:}
\cref{fig:dv:the-binomial-model:tree,fig:dv:the-binomial-model:conv,fig:dv:the-binomial-model:american}
et les valeurs numériques de l'\cref{ex:dv:the-binomial-model:whiteboard}.
\begin{tutsteps}
\item \textbf{Facteurs de pas.} Une fonction renvoie $(u,d,p)$ pour chaque arbre et
refuse un arbre comportant un arbitrage.
\omcode{code/firm/binomial/firm_binomial.py}{21}{46}{Pas et probabilité risque-neutre des trois arbres.}\label{lst:dv:the-binomial-model:params}
\item \textbf{Induction arrière.} Un tableau par tranche de temps, raccourci d'une unité
à chaque pas~; l'exercice américain tient en un \texttt{np.maximum} par tranche. L'argument
de dividende est utilisé au chapitre~5.
\omcode{code/firm/binomial/firm_binomial.py}{49}{70}{Induction arrière vectorisée, européenne ou américaine.}\label{lst:dv:the-binomial-model:price}
\item \textbf{Exécuter} \texttt{dv\_binomial.whiteboard()} pour les valeurs aux nœuds de la
\cref{fig:dv:the-binomial-model:tree}, puis \texttt{error(n, method)} pour
$n=10,\dots,200$ et \texttt{fig\_binomial.py}.
\end{tutsteps}
\textbf{Pour aller plus loin.} Fixer le strike à 101 au lieu de 100 et observer
le déplacement de la phase de l'oscillation~; puis pricer le put américain avec la moyenne des
\omterm{def:dv:the-binomial-model:crr}{arbres de Cox--Ross--Rubinstein} à $n$ pas et à $(n+1)$ pas, et comparer avec Leisen--Reimer.
\end{tutorial}

\section{Construction~: le pricer sur arbre}\label{bld:dv:the-binomial-model:binomial}

\begin{build}
\bfield{Rôle} Le premier pricer d'options américaines de la firme miniature, et la
référence contre laquelle est testé le moteur aux différences finies du chapitre~22. \bfield{Interface}
\texttt{params(method, spot, strike, t, r, q, vol, n) ->
(u, d, p)} avec method valant crr, jr ou lr~; \texttt{price(spot, strike, t, r, vol, n,
right, american, q, method, dividends)}~; \texttt{node\_values(\ldots)} pour les petits
arbres.
\bfield{Règles} Capitalisation continue~; $n$ impair pour Leisen--Reimer~; un arbitrage dans
les facteurs lève une erreur au lieu de renvoyer un nombre~; mémoire en $O(n)$.
\bfield{Tests d'acceptation} \texttt{code/firm/binomial/tests/}~: parité sur l'arbre~;
convergence des trois arbres vers la formule~; Leisen--Reimer à moins de $10^{-4}$ pour
101 pas~; call américain égal au call européen sans dividende~; put américain supérieur au put
européen et à sa valeur d'exercice.
\bfield{Pour aller plus loin} Un arbre trinomial (chapitre~22) et un pricer qui renvoie delta et
gamma à partir des deux premières tranches de temps, sans coût supplémentaire.
\end{build}

\begin{omsources}
\item J.~C.~Cox, S.~A.~Ross et M.~Rubinstein, ``Option pricing: a simplified
approach'', \emph{Journal of Financial Economics} 7 (1979) 229--263.
\item R.~J.~Rendleman et B.~J.~Bartter, ``Two-state option pricing'',
\emph{Journal of Finance} 34 (1979) 1093--1110.
\item R.~A.~Jarrow et A.~Rudd, \emph{Option Pricing}, Irwin, 1983.
\item D.~P.~J.~Leisen et M.~Reimer, ``Binomial models for option valuation:
examining and improving convergence'', \emph{Applied Mathematical Finance} 3
(1996) 319--346.
\end{omsources}

\section{Exercices}

\begin{exercise}[$\star$]\label{exo:dv:the-binomial-model:1}
Une action à 50 vaudra 60 ou 40~; les taux sont nuls. Donner le \omterm{def:dv:no-arbitrage-and-the-fundamental-theorems:claim}{portefeuille de réplication}
et le prix d'un call de strike 50, ainsi que la probabilité risque-neutre.
\end{exercise}

\begin{exercise}[$\star$]\label{exo:dv:the-binomial-model:2}
Donner $u$, $d$ et $p$ d'un pas de Cox--Ross--Rubinstein d'un mois avec $\sigma=20\%$
et $r=3\%$.
\end{exercise}

\begin{exercise}[$\star$]\label{exo:dv:the-binomial-model:3}
Vérifier la parité call-put sur l'arbre du tableau blanc.
\end{exercise}

\begin{exercise}[$\star\star$]\label{exo:dv:the-binomial-model:4}
Un collègue estime que l'action du tableau blanc monte avec une probabilité de
0.8, et non 0.6, et en conclut que le call vaut donc plus de 10.3603. Expliquer pourquoi
le prix ne dépend pas de cette estimation.
\end{exercise}

\begin{exercise}[$\star\star$]\label{exo:dv:the-binomial-model:5}
Calculer sur l'arbre du tableau blanc le prix d'une option digitale qui paie 1 si l'action termine au-dessus de 100.
\end{exercise}

\begin{exercise}[$\star\star$]\label{exo:dv:the-binomial-model:6}
Comparer les facteurs et les probabilités des \omterm{def:dv:the-binomial-model:crr}{arbres de Cox--Ross--Rubinstein} et de Jarrow--Rudd pour
$\sigma=20\%$, $r=5\%$, $\Delta t=0.1$.
\end{exercise}

\begin{exercise}[$\star\star\star$]\label{exo:dv:the-binomial-model:7}
\emph{Programmation.} Calculer le prix d'un put européen à la monnaie et d'échéance un an ($r=5\%$,
$\sigma=20\%$) avec 101 pas sur les \omterm{def:dv:the-binomial-model:crr}{arbres de Cox--Ross--Rubinstein} et de Leisen--Reimer
et comparer les deux valeurs à 5.5735.
\end{exercise}

\begin{exercise}[$\star\star\star$]\label{exo:dv:the-binomial-model:8}
\emph{Trouver le défaut.} «~Nous avons calibré l'arbre sur les données~: hausse de 5~\%, baisse de 3~\%, et
6~\% de croissance du cash par pas. La probabilité sort à 1.125, mais nous évaluons
quand même avec elle~».
\end{exercise}

\section{Problème~: l'arbre du tableau blanc}

\begin{problem}[{Problème du week-end --- de trois pas à un pricer opérationnel}]\label{pb:dv:the-binomial-model:1}
Un candidat reçoit l'arbre du tableau blanc de ce chapitre (action à 100, hausse de 10~\%,
baisse de 10~\%, croissance du cash de 2~\% par pas, trois pas, prix d'exercice 100), puis on lui demande de transformer
cet arbre en un pricer pour une option d'échéance un an avec $r=5\%$ et $\sigma=20\%$.

\medskip\noindent\textbf{Partie I --- Les options européennes.}
\begin{enumerate}
\item Donner la probabilité risque-neutre d'une hausse.
\item Lister les prix de l'action après trois pas.
\item Calculer le prix du call.
\item Donner le \omterm{def:dv:no-arbitrage-and-the-fundamental-theorems:claim}{portefeuille de réplication} au premier nœud.
\item Calculer le prix du put européen et vérifier la parité call-put.
\end{enumerate}

\medskip\noindent\textbf{Partie II --- Exercice américain.}
\begin{enumerate}[resume]
\item À quels nœuds l'exercice anticipé du put est-il optimal~?
\item Calculer le prix du put américain.
\item Donner la prime d'exercice anticipé.
\item Le call américain vaut-il plus que le call européen sur cet arbre~?
\item Pourquoi le détenteur au nœud 90 préfère-t-il exercer~?
\end{enumerate}

\medskip\noindent\textbf{Partie III --- Un arbre réel.}
\begin{enumerate}[resume]
\item Donner la valeur de Black--Scholes du put à la monnaie d'échéance un an.
\item Donner les erreurs de l'\omterm{def:dv:the-binomial-model:crr}{arbre de Cox--Ross--Rubinstein} pour 100 et 101 pas.
\item Donner l'erreur de l'\omterm{def:dv:the-binomial-model:lr}{arbre de Leisen--Reimer} pour 101 pas.
\item Calculer le prix du put américain avec un \omterm{def:dv:the-binomial-model:lr}{arbre de Leisen--Reimer} à 2\,001 pas, et sa prime.
\item Pourquoi l'erreur de l'\omterm{def:dv:the-binomial-model:crr}{arbre de Cox--Ross--Rubinstein} change-t-elle de signe selon la parité de $n$~?
\end{enumerate}

\medskip\noindent\textbf{Partie IV --- Jugement.}
\begin{enumerate}[resume]
\item Combien de pas chaque arbre faut-il pour une précision d'un centime~?
\item Comment le temps de calcul croît-il avec $n$, et qu'est-ce que cela implique pour un desk
réévaluant 10\,000 options américaines~?
\item Que suppose un arbre qu'un marché réel n'offre pas~?
\item Énoncer le \emph{résultat nommé}~: la prime d'exercice anticipé sur l'arbre
du tableau blanc, et le nombre de pas de l'\omterm{def:dv:the-binomial-model:crr}{arbre de Cox--Ross--Rubinstein} au-delà duquel le put reste
à moins d'un centime de la formule.
\item En une phrase~: pourquoi la probabilité de l'arbre n'est-elle pas une prévision~?
\end{enumerate}
\end{problem}

\section{Questions d'entretien}

\begin{interviewq}[$\star$ \iqroles{trader, chercheur}]\label{iq:dv:the-binomial-model:1}
Une action à 100 passe à 120 ou à 80 en une période, les taux sont nuls. Donnez le prix d'un call
de prix d'exercice 100 et indiquez la couverture.
\end{interviewq}

\begin{interviewq}[$\star$ \iqroles{chercheur}]\label{iq:dv:the-binomial-model:2}
Pourquoi la probabilité réelle d'une hausse n'intervient-elle pas dans le prix binomial~?
\end{interviewq}

\begin{interviewq}[$\star\star$ \iqroles{chercheur, développeur}]\label{iq:dv:the-binomial-model:3}
Pourquoi le prix de l'\omterm{def:dv:the-binomial-model:crr}{arbre de Cox--Ross--Rubinstein} oscille-t-il quand le nombre de pas augmente,
et comment le corrigez-vous~?
\end{interviewq}

\begin{interviewq}[$\star\star$ \iqroles{trader}]\label{iq:dv:the-binomial-model:4}
Pourquoi pourriez-vous exercer un put américain de façon anticipée, mais jamais un call américain sur
une action qui ne verse pas de dividende~?
\end{interviewq}

\begin{interviewq}[$\star\star$ \iqroles{développeur}]\label{iq:dv:the-binomial-model:5}
Combien de nœuds possède un \omterm{def:dv:the-binomial-model:recombining}{arbre recombinant} à $n$ pas, et combien de mémoire
le pricing sur cet arbre exige-t-il~?
\end{interviewq}

\begin{interviewq}[$\star\star\star$ \iqroles{développeur, chercheur}]\label{iq:dv:the-binomial-model:6}
Votre pricer sur arbre doit renvoyer le delta et le gamma en plus du prix, sans
réévaluer. Comment procédez-vous~?
\end{interviewq}
